\documentclass[11pt]{article}
\usepackage[margin=1in]{geometry}
\usepackage{amsmath,amssymb,amsthm,hyperref,listings}
\newtheorem{theorem}{Theorem}
\newtheorem{lemma}{Lemma}
\lstset{basicstyle=\ttfamily\footnotesize,columns=fullflexible,keepspaces=true,
        showstringspaces=false,breaklines=true}
\title{A variance proof of Davis's polynomial-ratio bound}
\author{Principia Math In Progress}
\date{17 September 2026}
\begin{document}
\maketitle
\noindent MathDB \href{https://mathdb.com/p/319187}{319187};
Davis, Conjecture 3.1 \cite{D}.

\section*{Task 1. Prove the original bound}
Write $(x)_{\underline i}=x(x-1)\cdots(x-i+1)$, with
$(x)_{\underline0}=1$, and let
\[
 P_d(x)=\sum_{i=0}^d
 \frac{(x)_{\underline i}^{\,2}(x)_{\underline{d-i}}^{\,2}}{i!(d-i)!},
 \qquad
 q_d(x)=\frac{2x^2-2dx}{d+1}+\frac d2.
\]
The remainder in the source is $R_d=P_{d+1}-q_dP_d$:
Lemma 1.1 of its journal version proves that $q_d$ is the polynomial
division quotient.

\begin{theorem}\label{thm:main}
For every integer $d\ge5$ and every real
$x\ge(d+\sqrt{d+2})/2$,
\begin{equation}\label{eq:target}
 -\frac{R_d(x)}{P_d(x)}>
       \frac{199(d-2)}{200(2x+d-2)}.
\end{equation}
Thus the conjectured weak inequality holds, with its exact constant
$0.995=199/200$.
\end{theorem}
The proof is analytic for all $d\ge40$. For the remaining 35 degrees,
we give exact integer polynomial certificates valid on whole real
rays, not pointwise numerical sampling.

\subsection*{The ratio as a variance}
Put $c=x-d/2$ and introduce the nonnegative weights
\[
 w_i=\frac{(x)_{\underline i}^{\,2}(x)_{\underline{d-i}}^{\,2}}{i!(d-i)!},
 \qquad \Pr(I=i)=\frac{w_i}{P_d(x)}\quad(0\le i\le d).
\]
The denominator is positive: the middle summand has no vanishing
factor when $x>d/2$. Since $w_i=w_{d-i}$, the variable $Z=I-d/2$
has mean zero. Write $V=\mathbb E Z^2$ and $M_4=\mathbb E Z^4$.

Adding a factor to either falling factorial gives
\[
 (d+1)P_{d+1}(x)
   =\sum_{i=0}^d w_i\bigl((x-i)^2+(x-d+i)^2\bigr).
\]
For example, the first sum contributes $j$ times the $j$th summand
of $P_{d+1}$, and the second contributes $d+1-j$ times that summand.
Consequently
\begin{equation}\label{eq:variance}
 \frac{P_{d+1}}{P_d}=\frac{2c^2+2V}{d+1},\qquad
 -\frac{R_d}{P_d}=\frac{d-4V}{2(d+1)}.
\end{equation}

There is also an exact fourth-moment identity. Define
\[
 a_i=(d-i)(x-i)^2,\qquad b_i=i(x-d+i)^2.
\]
Direct cancellation in the weights proves
$w_i a_i=w_{i+1}b_{i+1}$ for $0\le i<d$; this is a polynomial
identity and remains valid when some weights vanish. Since $a_d=b_0=0$,
telescoping gives
\[
 \mathbb E\{a_I(f(I+1)-f(I))+b_I(f(I-1)-f(I))\}=0.
\]
Use $f(i)=(i-d/2)^2$. In centered coordinates,
\[
 a_i-b_i=-2Z(c^2+dc+Z^2),\qquad
 a_i+b_i=dc^2+(d+4c)Z^2.
\]
It follows that
\begin{equation}\label{eq:moment}
 4M_4+A V=dc^2,\qquad A=4c^2+4(d-1)c-d.
\end{equation}
Cauchy--Schwarz gives $M_4\ge V^2$. Thus, with $t=4V$,
\begin{equation}\label{eq:quadratic}
 t^2+At\le4dc^2.
\end{equation}

\subsection*{Every degree at least 40}
Suppose $d\ge40$. The required domain implies
$c\ge\sqrt{d+2}/2>3$, so $A>0$. Set
\[
 T=d-\frac{199(d+1)(d-2)}{200(c+d-1)}.
\]
Here $T>d-(d+1)(d-2)/(d+2)=3-4/(d+2)>0$.
We will prove $T^2+AT>4dc^2$. Since $y\mapsto y^2+Ay$ is strictly
increasing for $y\ge0$, \eqref{eq:quadratic} will then imply $t<T$.

Put $u=d-40\ge0$ and $v=c-3\ge0$. An exact expansion is
\begin{equation}\label{eq:largecert}
 40000(c+d-1)^2(T^2+AT-4dc^2)
       =C_0(u)+C_1(u)v+C_2(u)v^2+C_3(u)v^3,
\end{equation}
where
\begin{align*}
 C_0(u)&=2201u^4+240358u^3+8708957u^2
                  +110582964u+241318564,\\
 C_1(u)&=800u^4+95400u^3+4131000u^2
                  +80610800u+677526400,\\
 C_2(u)&=1600u^3+196000u^2+8631200u+136276800,\\
 C_3(u)&=800u^2+63200u+1566400.
\end{align*}
Every displayed coefficient is positive. Hence \eqref{eq:largecert}
is strictly positive, proving $4V=t<T$. By \eqref{eq:variance},
\[
 -\frac{R_d}{P_d}>\frac{199(d-2)}{400(c+d-1)}
                =\frac{199(d-2)}{200(2x+d-2)}.
\]
In fact this part of the proof works whenever $d\ge40$ and $c\ge3$.

\subsection*{Exact certificates for degrees 5 through 39}
Fix such a degree and put
\[
 b=\lfloor d/2\rfloor,\quad q=\lceil d/2\rceil,\quad
 K=1000d+\left\lfloor\sqrt{10^6(d+2)}\right\rfloor,
 \quad X=2000x-K.
\]
The source domain implies $X\ge0$. Define the degree-$b$ integer
polynomial
\[
 A_i(X)=
 \prod_{j=0}^{\min(i,d-i)-1}(X+K-2000j)
 \prod_{j=q}^{\max(i,d-i)-1}(X+K-2000j),
\]
where an empty product is one, and set
\[
 S(X)=\sum_{i=0}^d\binom di A_i(X)^2,\qquad
 N(X)=\sum_{i=0}^d\bigl(d-(2i-d)^2\bigr)\binom di A_i(X)^2.
\]
Factoring the common square $(x)_{\underline q}^{\,2}$ from the
weights shows that $S>0$ and $N/S=d-4V$. Thus
\eqref{eq:target} is equivalent to positivity of
\begin{equation}\label{eq:smallcert}
 E_d(X)=(X+K+1000(d-2))N(X)-1990(d+1)(d-2)S(X).
\end{equation}
For each of the 35 indicated degrees, every coefficient of $E_d$ is
nonnegative and its constant coefficient is positive. The degree is
$2\lfloor d/2\rfloor$; there are 787 coefficients in total.
These finite integer assertions are verified by the complete checker
below. The coefficient vectors, divided by their positive integer
contents, are also recorded in \texttt{polynomial-certificates.json}.

\begin{lstlisting}[language=Python]
from math import comb, isqrt

def mul(p, q):
    r = [0] * (len(p) + len(q) - 1)
    for i, a in enumerate(p):
        for j, b in enumerate(q):
            r[i+j] += a*b
    return r

total = 0
for d in range(5, 40):
    b, q = d//2, (d+1)//2
    K = 1000*d + isqrt(1000000*(d+2))
    S, N = [0]*(2*b+1), [0]*(2*b+1)
    for i in range(d+1):
        p = [1]
        J = list(range(min(i,d-i))) + list(range(q,max(i,d-i)))
        for j in J:
            p = mul(p, [K-2000*j, 1])
        for j, a in enumerate(mul(p,p)):
            w = comb(d,i)*a
            S[j] += w
            N[j] += (d-(2*i-d)**2)*w
    E = mul(N, [K+1000*(d-2), 1])
    for j, a in enumerate(S):
        E[j] -= 1990*(d+1)*(d-2)*a
    while len(E)>1 and E[-1] == 0:
        E.pop()
    assert len(E)-1 == 2*b
    assert E[0] > 0 and all(a >= 0 for a in E)
    total += len(E)
assert total == 787
\end{lstlisting}
This code uses only exact integer arithmetic from Python's standard
library. Its coefficient assertions imply $E_d(X)>0$ for every real
$X\ge0$, and therefore prove \eqref{eq:target} for every real argument
in the required domain. Together with the large-degree argument,
this completes the proof.\hfill$\square$

\paragraph{Scope and verification.}
The proof handles every integer degree $d\ge5$, includes the stated
cutoff, and proves a strict inequality. The accompanying verifier
checks the universal symbolic identities, the 17 positive terms in
\eqref{eq:largecert}, all 35 small-degree certificates, and 70 exact
rational evaluations from the original falling-factorial definition.
Those last evaluations check normalization only; the polynomial
certificates prove the full real intervals. The same agent authored
and audited the argument; independent peer review is not claimed.

\section*{Task 2. Determine the optimal universal constant}
Determine the infimum, over the full source domain, of
\[
 \frac{2x+d-2}{d-2}\left(-\frac{R_d(x)}{P_d(x)}\right).
\]
In particular, prove or refute the candidate that it is attained at
$d=27$ and $x=(27+\sqrt{29})/2$. Obtain an exact algebraic
description and a certified interval for the optimal constant.

\section*{Task 3. Remove the finite coefficient check}
Use additional detailed-balance moment identities or a sharper lower
bound for $\mathbb E Z^4$ to prove the stated $199/200$ inequality
uniformly from $d=5$ onward. Seek a short analytic certificate that
replaces the 35 polynomial expansions while preserving the full
real-argument domain.

\section*{Task 4. Prove a uniform two-sided approximation}
Derive explicit upper and lower error bounds for
$-R_d/P_d-(d-2)/(2x+d-2)$, uniform for
$x\ge(d+\sqrt{d+2})/2$. Determine the leading correction both when
$x-d/2$ is proportional to $\sqrt d$ and when $x/d$ stays above
$1/2$, with rigorous control through the transition between them.

\section*{Task 5. Resolve the extremum conjecture}
For odd $d\ge5$, determine the number of real critical points of
\[
 Q_d(x)=1+\frac{R_d(x)}{P_d(x)}\frac{2x+d-2}{d-2},
\]
interpreted after cancelling common factors. Prove or refute the
source's single-maximum, single-minimum assertion and derive rigorous
large-$d$ expansions for their locations. This is Conjecture 3.2 in
the journal version, numbered 3.9 in the preprint.

\begin{thebibliography}{9}
\bibitem{D}
D. M. Davis, \emph{Maximizing a combinatorial expression arising from
crowd estimation}, International Journal of Modern Mathematics
\textbf{5} (2010), 129--148. Author-hosted journal version, revised
14 May 2010:
\href{https://www.lehigh.edu/~dmd1/fredijmm.pdf}{\texttt{fredijmm.pdf}},
Lemma 1.1 on PDF p. 5 and Conjecture 3.1 on PDF p. 12.
Preprint: \href{https://arxiv.org/abs/0910.5725}{arXiv:0910.5725v1}.
\end{thebibliography}
\end{document}
